% ============================================================
%  Chapter 8 — Span and Generating Systems
% ============================================================
\section{Span and Generating Systems}
\label{sec:span}

Given a vector space $V$ over $\F$, we often want to describe a subspace by
specifying a small list of vectors that ``generate'' it.  The notion of a
\emph{linear combination} formalises what it means for one vector to be built
from others, and the \emph{span} of a set is the collection of all such
combinations.

% ---- Linear combination ------------------------------------
\begin{defbox}[Definition --- Linear Combination]
\label{def:linear-combination}
Let $V$ be a vector space over $\F$ and let $v_1,\ldots,v_n\in V$.  A vector
$v\in V$ is a \textbf{linear combination} of $v_1,\ldots,v_n$ if there exist
scalars $\alpha_1,\ldots,\alpha_n\in\F$ such that
\[
    v = \alpha_1 v_1 + \alpha_2 v_2 + \cdots + \alpha_n v_n
      = \sum_{i=1}^{n}\alpha_i v_i.
\]
\end{defbox}

% ---- Span --------------------------------------------------
\begin{defbox}[Definition --- Span of Vectors]
\label{def:span}
The \textbf{span} of $v_1,\ldots,v_n\in V$ is the set of all their linear
combinations:
\[
    \spn\{v_1,\ldots,v_n\}
    = \left\{\, \sum_{i=1}^{n}\alpha_i v_i \;\middle|\; \alpha_1,\ldots,\alpha_n\in\F \,\right\}.
\]
The set $\{v_1,\ldots,v_n\}$ is called a \textbf{generating system}
(or \textbf{spanning set}) of $\spn\{v_1,\ldots,v_n\}$.
More generally, for any subset $S\subseteq V$, $\spn(S)$ is the set of all
finite linear combinations of elements of $S$.
\end{defbox}

% ---- Span is a subspace ------------------------------------
\begin{propbox}[Proposition --- Span is a Subspace]
\label{prop:span-subspace}
For any $v_1,\ldots,v_n\in V$, the set $\spn\{v_1,\ldots,v_n\}$ is a subspace
of $V$.  In fact it is the \emph{smallest} subspace of $V$ containing every
$v_i$.
\end{propbox}

\begin{proofbox}[Proof --- Span is a Subspace]
\begin{enumerate}[label=(\roman*)]
    \item \textbf{Zero vector:}
          $\mathbf{0} = 0\cdot v_1 + \cdots + 0\cdot v_n \in \spn\{v_1,\ldots,v_n\}$.
    \item \textbf{Closed under addition:}
          If $u = \sum\alpha_i v_i$ and $w = \sum\beta_i v_i$, then
          $u+w = \sum(\alpha_i+\beta_i)v_i \in \spn\{v_1,\ldots,v_n\}$.
    \item \textbf{Closed under scalar multiplication:}
          $\lambda u = \sum(\lambda\alpha_i)v_i \in \spn\{v_1,\ldots,v_n\}$.
\end{enumerate}
Any subspace containing all $v_i$ must contain every linear combination of
them, so $\spn\{v_1,\ldots,v_n\}$ is minimal.
\hfill$\blacksquare$
\end{proofbox}

% ---- Column-space / matrix view ----------------------------
\begin{tipbox}[Remark --- Span in $\F^m$ via Matrices]
\label{rem:span-matrix}
When $V = \F^m$ and $v_1,\ldots,v_n\in\F^m$, form the matrix
$A = [v_1\mid\cdots\mid v_n]\in\Mmn{m}{n}$.  Then
\[
    \spn\{v_1,\ldots,v_n\}
    = \{b\in\F^m \mid Ax = b \text{ is consistent}\}.
\]
That is, the span equals the column space of $A$.
\end{tipbox}

% ---- Standard examples -------------------------------------
\begin{exbox}[Example --- Span of Standard Vectors]
\begin{enumerate}
    \item In $\R^2$: $\spn\{(1,0),(0,1)\} = \R^2$, since
          $(x,y) = x(1,0)+y(0,1)$.

    \item In $\F^n$: $\spn\{e_1,\ldots,e_n\} = \F^n$ where $e_i$ are standard
          basis vectors.

    \item In $\Mmn{m}{n}$: let $E^{ij}$ be the matrix with a $1$ in position
          $(i,j)$ and $0$ elsewhere.  Then
          $\spn\{E^{ij}\mid 1\le i\le m,\,1\le j\le n\} = \Mmn{m}{n}$.

    \item In $\F[x]$: $\F[x] = \spn\{x^i\mid i\in\N_0\}$, since every polynomial
          is a finite linear combination of monomials.
\end{enumerate}
\end{exbox}

% ---- Intersection of subspaces -----------------------------
\begin{propbox}[Proposition --- Intersection of Subspaces]
\label{prop:intersection-subspace}
If $S$ and $T$ are subspaces of $V$, then $S\cap T$ is also a subspace of $V$.
\end{propbox}

\begin{proofbox}[Proof --- Intersection of Subspaces]
Since $\mathbf{0}\in S$ and $\mathbf{0}\in T$, we have $\mathbf{0}\in S\cap T$.
If $u,v\in S\cap T$ and $\alpha\in\F$, then $u+v\in S$ (as $S\le V$) and
$u+v\in T$ (as $T\le V$), so $u+v\in S\cap T$; similarly $\alpha v\in S\cap T$.
\hfill$\blacksquare$
\end{proofbox}

% ---- Union is NOT a subspace -------------------------------
\begin{warnbox}[Warning --- Union of Subspaces Need Not Be a Subspace]
In general, $S\cup T$ is \emph{not} a subspace of $V$.

\textbf{Counterexample.}  Let $V=\R^2$,\; $S=\{(x,0)\mid x\in\R\}$,\;
$T=\{(0,y)\mid y\in\R\}$.  Both are subspaces, but
$(1,0)\in S\cup T$ and $(0,1)\in S\cup T$ while
$(1,0)+(0,1)=(1,1)\notin S\cup T$.
\end{warnbox}

% ---- Rank and generating sets in F^m -----------------------
\begin{propbox}[Proposition --- Rank Criterion for Generating Sets]
\label{prop:rank-generating}
Vectors $v_1,\ldots,v_n\in\F^m$ form a generating set for $\F^m$ if and only
if the matrix $A=[v_1\mid\cdots\mid v_n]$ has rank $m$.  In particular,
$m\le n$ must hold whenever such a generating set exists.
\end{propbox}

\begin{proofbox}[Proof --- Rank Criterion for Generating Sets]
The vectors span $\F^m$ iff every $b\in\F^m$ is a linear combination of the
$v_i$, i.e.\ iff $Ax=b$ is consistent for every $b$.  By the theory of linear
systems this is equivalent to the RREF of $A$ having a pivot in every row,
i.e.\ $\rk(A)=m$.  Since $A$ has $n$ columns, $\rk(A)\le n$, so $m\le n$.
\hfill$\blacksquare$
\end{proofbox}

% ---- Subset span relation ----------------------------------
\begin{tipbox}[Remark --- Span and Subsets]
If $w_1,\ldots,w_k\in\spn\{v_1,\ldots,v_n\}$, then
$\spn\{w_1,\ldots,w_k\}\subseteq\spn\{v_1,\ldots,v_n\}$.
Consequently, if $v_r\in\spn\{v_1,\ldots,v_{r-1}\}$, then
$\spn\{v_1,\ldots,v_r\}=\spn\{v_1,\ldots,v_{r-1}\}$ (the redundant vector
can be removed without changing the span).
\end{tipbox}
